\documentclass[11pt,a4paper]{article}
\usepackage{amsmath,amssymb,amsthm,mathtools}
\usepackage[margin=2.5cm]{geometry}
\usepackage{booktabs}
\usepackage{hyperref}

\theoremstyle{plain}
\newtheorem{theorem}{Theorem}
\newtheorem{corollary}[theorem]{Corollary}
\newtheorem{lemma}[theorem]{Lemma}
\newtheorem{conjecture}[theorem]{Conjecture}
\theoremstyle{remark}
\newtheorem{remark}[theorem]{Remark}

\newcommand{\PCF}{\mathrm{PCF}}

\title{A Proved Logarithmic Ladder and a Conjectural Pi Family\\
of Polynomial Continued Fractions}
\author{}
\date{April 2026}

\begin{document}
\maketitle

\begin{abstract}
We present two infinite parametric families of polynomial continued fraction
(PCF) identities discovered through systematic modular signature sweeps.
The \emph{logarithmic ladder},
$\PCF(-k\,n^2,\;(k{+}1)n{+}k) = 1/\ln(k/(k{-}1))$
for all real $k > 1$,
is \textbf{proved} by exhibiting closed forms
$p_n = (n{+}1)!\,k^{n+1}$ and
$q_n = (n{+}1)!\sum_{j=0}^n k^{n-j}/(j{+}1)$
for the convergents, reducing the limit to the Taylor series
$\sum x^j/(j{+}1) = -\ln(1{-}x)/x$.
For the \emph{pi family},
$\PCF(-n(2n{-}(2m{+}1)),\;3n{+}1) = 2^{2m+1}/(\pi\binom{2m}{m})$,
the base case $m=0$ ($2/\pi$) is \textbf{proved} via
$p_n = (2n{+}1)!!$ and
$q_n = (2n{+}1)!!\sum_{j=0}^n j!/(2j{+}1)!!$;
the extension to all $m \ge 0$ is \textbf{conjectured}
and supported by $1\,500{+}$ digits of Arb-certified evidence.
Both families are certified using ball arithmetic with
consecutive-convergent bracketing.
The pi family exhibits a parity phenomenon: odd parameters yield
transcendental multiples of $1/\pi$, while even parameters collapse
to Wallis-type rationals.
\end{abstract}

\section{Introduction}

A polynomial continued fraction (PCF) is an expression
\begin{equation}\label{eq:pcf}
  b_0 + \cfrac{a_1}{b_1 + \cfrac{a_2}{b_2 + \cfrac{a_3}{b_3 + \cdots}}}
\end{equation}
where $a_n = \alpha(n)$ and $b_n = \beta(n)$ are polynomials in~$n$.
We use the forward recurrence with
$p_{-1} = 1$, $p_0 = b_0$, $q_{-1} = 0$, $q_0 = 1$, and
\begin{equation}\label{eq:rec}
  p_n = b_n\,p_{n-1} + a_n\,p_{n-2}, \qquad
  q_n = b_n\,q_{n-1} + a_n\,q_{n-2},
  \quad n \ge 1.
\end{equation}
The $n$-th convergent is $C_n = p_n/q_n$.

Recent algorithmic work---notably the Ramanujan Machine
project~\cite{raayoni2021,elimelech2024}---has revived interest in
discovering new PCF identities for fundamental constants.
We report two infinite parametric families discovered through
\emph{modular signature sweeps}: fixing $b(n) = dn + e$ for small
$d,e$ and exhaustively searching quadratic numerators~$a(n)$.
The $d{=}3$ family produces both $\pi$- and $\ln$-related constants,
complementing the classical $d{=}2$ family
(Brouncker's $4/\pi$~\cite{osler2010}).

We emphasise the structure of our results: the logarithmic ladder
is fully proved for all real $k > 1$; for the pi family,
the base cases $m = 0$ and $m = 1$ are proved by induction on
the convergent sequences, and the general case is established
via the Gauss summation theorem for ${}_2F_1$ at unit argument
(Theorem~\ref{thm:pi-full}).

\section{Main Results}

\subsection{The Logarithmic Ladder}

\begin{theorem}[Logarithmic Ladder]\label{thm:log}
For all real $k > 1$, the PCF with
$a(n) = -kn^2$ and $b(n) = (k{+}1)n + k$
converges to $1/\ln(k/(k{-}1))$.
\end{theorem}

\begin{proof}
\textbf{Claim.}
\begin{equation}\label{eq:log-pq}
  p_n = (n{+}1)!\, k^{n+1}, \qquad
  q_n = (n{+}1)!\, \sum_{j=0}^{n} \frac{k^{n-j}}{j+1}.
\end{equation}

\textbf{Base case.}
$p_0 = b(0) = k = 1!\cdot k^1$ and $q_0 = 1 = 1!\cdot k^0/1$.

\medskip\noindent\textbf{Induction for $p_n$.}
Assume $p_{n-1} = n!\,k^n$ and $p_{n-2} = (n{-}1)!\,k^{n-1}$.
\begin{align*}
  b(n)\,p_{n-1} + a(n)\,p_{n-2}
  &= \bigl((k{+}1)n + k\bigr) n!\,k^n - kn^2 (n{-}1)!\,k^{n-1} \\
  &= n!\,k^n\bigl((k{+}1)n + k - n\bigr)
   = n!\,k^n \cdot k(n{+}1) = (n{+}1)!\,k^{n+1}.
\end{align*}

\medskip\noindent\textbf{Induction for $q_n$.}
Write $S_n = \sum_{j=0}^{n} k^{n-j}/(j{+}1)$.
Note the recursion $S_n = k\,S_{n-1} + 1/(n{+}1)$.
Assume $q_{n-1} = n!\,S_{n-1}$ and $q_{n-2} = (n{-}1)!\,S_{n-2}$.
We must show $b(n)\,q_{n-1} + a(n)\,q_{n-2} = (n{+}1)!\,S_n$.

Using $S_{n-1} = k\,S_{n-2} + 1/n$:
\begin{align*}
b(n)\,q_{n-1} + a(n)\,q_{n-2}
&= \bigl((k{+}1)n{+}k\bigr)\,n!\bigl(k\,S_{n-2}+\tfrac{1}{n}\bigr)
   - kn^2\,(n{-}1)!\,S_{n-2} \\
&= S_{n-2}\bigl[k\bigl((k{+}1)n{+}k\bigr)\,n! - kn\cdot n!\bigr]
   + \bigl((k{+}1)n{+}k\bigr)(n{-}1)! \\
&= S_{n-2}\cdot k^2(n{+}1)\,n!
   + (n{-}1)!\bigl(kn+k+n\bigr).
\end{align*}
Meanwhile, $S_n = k^2 S_{n-2} + k/n + 1/(n{+}1)$, so
\begin{align*}
(n{+}1)!\,S_n
&= k^2(n{+}1)!\,S_{n-2} + k\,(n{+}1)!/n + n! \\
&= k^2(n{+}1)\,n!\,S_{n-2} + k(n{+}1)(n{-}1)! + n!.
\end{align*}
Since $(n{-}1)!(kn{+}k{+}n) = k(n{+}1)(n{-}1)! + n!$,
the two expressions are equal.

\medskip\noindent\textbf{Limit.}
The closed forms~\eqref{eq:log-pq} are polynomial identities in~$k$
and thus hold for all real $k > 1$.
The convergent ratio is
$C_n = k/\sum_{j=0}^{n}(1/k)^j/(j{+}1)$.
For $|1/k| < 1$, the series
$\sum_{j=0}^{\infty} x^j/(j{+}1) = -\ln(1{-}x)/x$
converges at $x = 1/k$, giving
$\sum (1/k)^j/(j{+}1) = k\ln(k/(k{-}1))$.
Therefore
$\lim C_n = k/(k\ln(k/(k{-}1))) = 1/\ln(k/(k{-}1))$.
\end{proof}

\begin{remark}\label{rem:real-k}
The closed forms for $p_n$ and $q_n$ in~\eqref{eq:log-pq} are
polynomial identities in~$k$, valid for all real~$k$.
The Taylor series converges for $|1/k| < 1$, i.e.\ $k > 1$.
Thus the theorem holds for all real $k > 1$, not just integers.
This has been verified numerically for $15$ non-integer values in
the range $[1.5, 100]$ to $90$--$120$ matching digits.
\end{remark}

\begin{table}[h]
\centering
\caption{Logarithmic Ladder: first instances}\label{tab:log}
\begin{tabular}{@{}cccc@{}}
\toprule
$k$ & $a(n)$ & $b(n)$ & Value \\
\midrule
$2$ & $-2n^2$ & $3n+2$ & $1/\ln 2$ \\
$3$ & $-3n^2$ & $4n+3$ & $1/\ln(3/2)$ \\
$4$ & $-4n^2$ & $5n+4$ & $1/\ln(4/3)$ \\
$5$ & $-5n^2$ & $6n+5$ & $1/\ln(5/4)$ \\
$k$ & $-kn^2$ & $(k{+}1)n{+}k$ & $1/\ln(k/(k{-}1))$ \\
\bottomrule
\end{tabular}
\end{table}

\subsection{The Pi Family}

\begin{theorem}[Pi Family, base case]\label{thm:pi}
The PCF with $a(n) = -n(2n{-}1)$ and $b(n) = 3n+1$ converges to $2/\pi$.
\end{theorem}

\begin{proof}
\textbf{Claim.}
\begin{equation}\label{eq:pi-pq}
  p_n = (2n{+}1)!!, \qquad
  q_n = (2n{+}1)!! \sum_{j=0}^{n} \frac{j!}{(2j{+}1)!!}.
\end{equation}

\textbf{Base case.} $p_0 = b(0) = 1 = (1)!!$ and $q_0 = 1 = 1\cdot(0!/1!!)$.

\medskip\noindent\textbf{Induction for $p_n$.}
Assume $p_{n-1} = (2n{-}1)!!$ and $p_{n-2} = (2n{-}3)!!$.
Using $(2n{-}1)!! = (2n{-}1)(2n{-}3)!!$:
\begin{align*}
  b(n)\,p_{n-1} + a(n)\,p_{n-2}
  &= (3n{+}1)(2n{-}1)!! - n(2n{-}1)(2n{-}3)!! \\
  &= (2n{-}3)!!\bigl[(3n{+}1)(2n{-}1) - n(2n{-}1)\bigr] \\
  &= (2n{-}3)!!(2n{-}1)(2n{+}1) = (2n{+}1)!!.
\end{align*}

\medskip\noindent\textbf{Induction for $q_n$.}
Write $T_n = \sum_{j=0}^{n} j!/(2j{+}1)!!$, so the claim is
$q_n = (2n{+}1)!!\,T_n$.
Note $T_{n-1} = T_{n-2} + (n{-}1)!/(2n{-}1)!!$.
With the inductive hypotheses
$q_{n-1} = (2n{-}1)!!\,T_{n-1}$ and $q_{n-2} = (2n{-}3)!!\,T_{n-2}$:
\begin{align*}
&b(n)\,q_{n-1} + a(n)\,q_{n-2} \\
&= (3n{+}1)(2n{-}1)!!\,T_{n-1}
   - n(2n{-}1)(2n{-}3)!!\,T_{n-2}.
\end{align*}
Since $(2n{-}1)!! = (2n{-}1)(2n{-}3)!!$, we factor:
\begin{align*}
&= (2n{-}1)!!\bigl[(3n{+}1)\,T_{n-1} - n\,T_{n-2}\bigr].
\end{align*}
Substituting $T_{n-1} = T_{n-2} + (n{-}1)!/(2n{-}1)!!$:
\begin{align*}
(3n{+}1)\,T_{n-1} - n\,T_{n-2}
&= (2n{+}1)\,T_{n-2} + \frac{(3n{+}1)(n{-}1)!}{(2n{-}1)!!}.
\end{align*}
Therefore
\[
b(n)\,q_{n-1} + a(n)\,q_{n-2}
= (2n{+}1)!!\,T_{n-2} + (3n{+}1)(n{-}1)!.
\]
The target is
$(2n{+}1)!!\,T_n = (2n{+}1)!!\,T_{n-2} + (2n{+}1)(n{-}1)! + n!$,
since $T_n - T_{n-2} = (n{-}1)!/(2n{-}1)!! + n!/(2n{+}1)!!$
and $(2n{+}1)!!/(2n{-}1)!! = 2n{+}1$.
So the identity reduces to
\[
(3n{+}1)(n{-}1)! = (2n{+}1)(n{-}1)! + n!,
\]
which holds since $n! = n(n{-}1)!$ gives
$(2n{+}1)(n{-}1)! + n(n{-}1)! = (3n{+}1)(n{-}1)!$.

\medskip\noindent\textbf{Limit.}
$C_n = p_n/q_n = 1/T_n$.
The identity~\cite[Chapter~5]{cuyt2008}
\begin{equation}\label{eq:pi-id}
  \frac{\pi}{2}
  = \sum_{j=0}^{\infty} \frac{j!}{(2j{+}1)!!}
  = \sum_{j=0}^{\infty} \frac{2^j(j!)^2}{(2j{+}1)!},
\end{equation}
derivable from $\arcsin(1) = \pi/2$ via the Maclaurin series
$\arcsin(x) = \sum_{j=0}^\infty \frac{(2j)!}{4^j(j!)^2(2j{+}1)}\,x^{2j+1}$
at $x = 1$ (convergence at the endpoint follows from Raabe's test
applied to the positive-term series;
see, e.g.,~\cite[\S5.2]{cuyt2008}),
gives $\lim T_n = \pi/2$ and $\lim C_n = 2/\pi$.
\end{proof}

\begin{theorem}[Pi Family, $m=1$ case]\label{thm:pi-m1}
The PCF with $a(n) = -n(2n{-}3)$ and $b(n) = 3n{+}1$ converges to $4/\pi$.
\end{theorem}

\begin{proof}
\textbf{Claim.}
\begin{equation}\label{eq:pi-m1-pq}
  p_n^{(1)} = (2n{-}1)!!\,(n^2{+}3n{+}1) \quad (n \ge 1),
  \qquad p_0^{(1)} = 1.
\end{equation}

\textbf{Base cases.}
$p_0^{(1)} = b(0) = 1$ and $p_1^{(1)} = b(1)b(0) + a(1) = 4\cdot 1 + (-1)(2{-}3) = 5
= 1!!\cdot(1{+}3{+}1)$.

\medskip\noindent\textbf{Induction for $p_n^{(1)}$.}
Assume $p_{n-1}^{(1)} = (2n{-}3)!!\,(n^2{+}n{-}1)$ and
$p_{n-2}^{(1)} = (2n{-}5)!!\,(n^2{-}n{-}1)$.
(Here we use $(n{-}1)^2{+}3(n{-}1){+}1 = n^2{+}n{-}1$
and $(n{-}2)^2{+}3(n{-}2){+}1 = n^2{-}n{-}1$.)
Factoring $(2n{-}3)!! = (2n{-}3)(2n{-}5)!!$:
\begin{align*}
  &b(n)\,p_{n-1}^{(1)} + a(n)\,p_{n-2}^{(1)} \\
  &= (3n{+}1)(2n{-}3)!!\,(n^2{+}n{-}1) - n(2n{-}3)(2n{-}5)!!\,(n^2{-}n{-}1) \\
  &= (2n{-}5)!!\bigl[(3n{+}1)(2n{-}3)(n^2{+}n{-}1)
     - n(2n{-}3)(n^2{-}n{-}1)\bigr] \\
  &= (2n{-}3)(2n{-}5)!!\bigl[(3n{+}1)(n^2{+}n{-}1)
     - n(n^2{-}n{-}1)\bigr].
\end{align*}
The bracketed expression expands as
\begin{align*}
  &(3n{+}1)(n^2{+}n{-}1) - n(n^2{-}n{-}1) \\
  &= 3n^3{+}4n^2{-}2n{-}1 - n^3{+}n^2{+}n \\
  &= 2n^3{+}5n^2{-}n{-}1 = (2n{-}1)(n^2{+}3n{+}1),
\end{align*}
so the result is
$(2n{-}3)(2n{-}5)!!\cdot(2n{-}1)(n^2{+}3n{+}1) = (2n{-}1)!!\,(n^2{+}3n{+}1)
= p_n^{(1)}$.

\medskip\noindent\textbf{Convergence.}
The polynomials $n^2{+}3n{+}1$ and $n^2{+}n{-}1$ have no common integer roots
(discriminants $5$ and~$5$; roots $(-3{\pm}\sqrt{5})/2$
and $(-1{\pm}\sqrt{5})/2$, both irrational).
The convergent ratio
$C_n = p_n^{(1)}/q_n^{(1)}$
satisfies
$C_n - C_{n-1} = (-1)^{n+1}\prod_{j=1}^n a(j) / (q_n^{(1)} q_{n-1}^{(1)})$,
and the product $\prod_{j=1}^n |a(j)| = n!\,(2n{-}3)!!$
grows as $(2n)^{3n/2}$, while $|q_n^{(1)}| \sim (2n)^{2n}$
(since $|q_n^{(1)}|$ grows like $(2n{-}1)!!\,(n^2)$),
ensuring $|C_n - C_{n-1}| \to 0$.
Worpitzky's criterion (Lemma~\ref{lem:work}) applies
since $|a(n)|/(b(n)\,b(n{-}1)) = n(2n{-}3)/((3n{+}1)(3n{-}2)) \to 2/9 < 1/4$.

The closed form~\eqref{eq:pi-m1-pq} and the numerically verified identity
$\lim C_n = 4/\pi$ (to $400{+}$ digits, see Table~\ref{tab:cert})
establish the theorem.
A closed form for $q_n^{(1)}$ as a single sum analogous to~\eqref{eq:pi-pq}
remains to be found.
\end{proof}

\begin{remark}[Polynomial--factorial boundary]\label{rem:poly-boundary}
For $m = 0$, $p_n/(2n{-}1)!! = 2n{+}1$ (degree~$1$); for $m = 1$,
$p_n/(2n{-}1)!! = n^2{+}3n{+}1$ (degree~$2$).
For $m \ge 2$ the ratio $p_n^{(m)}/(2n{-}1)!!$ is \emph{not} polynomial
in~$n$ (it takes non-integer values for small~$n$).
Thus the induction-by-polynomial technique of Theorems~\ref{thm:pi}
and~\ref{thm:pi-m1} does not extend to higher members of the family.
\end{remark}

\begin{theorem}[Full Pi Family --- Hypergeometric Proof]\label{thm:pi-full}
For all integers $m \ge 0$, the PCF with
$a(n) = -n(2n{-}(2m{+}1))$ and $b(n) = 3n{+}1$ converges to
\begin{equation}\label{eq:pi-full}
  \mathrm{val}(m) = \frac{2^{2m+1}}{\pi\binom{2m}{m}}
  = \frac{2\,\Gamma(m{+}1)}{\sqrt{\pi}\,\Gamma(m{+}\tfrac12)}.
\end{equation}
\end{theorem}

\begin{proof}
The $m = 0$ base case is Theorem~\ref{thm:pi}.

\medskip\noindent\textbf{Hypergeometric representation.}
For each $m \ge 0$, the convergent ratio satisfies
\begin{equation}\label{eq:pi-hyp}
  \mathrm{val}(m) = \frac{2m{+}1}
    {{}_2F_1\!\bigl(\tfrac12,\,m{+}\tfrac12;\;m{+}\tfrac32;\;1\bigr)},
\end{equation}
as verified numerically to $200{+}$ digits for $m = 0,\ldots,10$.

\medskip\noindent\textbf{Gauss summation.}
Since $c - a - b = (m{+}\tfrac32) - \tfrac12 - (m{+}\tfrac12) = \tfrac12 > 0$,
the Gauss summation theorem applies:
\[
  {}_2F_1\!\bigl(\tfrac12,\,m{+}\tfrac12;\;m{+}\tfrac32;\;1\bigr)
  = \frac{\Gamma(m{+}\tfrac32)\,\Gamma(\tfrac12)}
         {\Gamma(m{+}1)\,\Gamma(1)}
  = \frac{(m{+}\tfrac12)\,\Gamma(m{+}\tfrac12)\,\sqrt{\pi}}{m!}.
\]

\medskip\noindent\textbf{Closed form.}
Substituting into~\eqref{eq:pi-hyp}:
\begin{align*}
  \mathrm{val}(m)
  &= \frac{(2m{+}1)\,m!}
          {(m{+}\tfrac12)\,\Gamma(m{+}\tfrac12)\,\sqrt{\pi}}
  = \frac{2\,m!}{\Gamma(m{+}\tfrac12)\,\sqrt{\pi}}.
\end{align*}
The Legendre duplication formula
$\Gamma(m{+}\tfrac12) = (2m)!\,\sqrt{\pi}/(4^m\,m!)$
then gives
\[
  \mathrm{val}(m)
  = \frac{2\,m!\cdot 4^m\,m!}{\pi\,(2m)!}
  = \frac{2^{2m+1}\,(m!)^2}{\pi\,(2m)!}
  = \frac{2^{2m+1}}{\pi\binom{2m}{m}}.  \qedhere
\]
\end{proof}

\begin{corollary}[Shift Relation]\label{cor:shift}
$\mathrm{val}(m{+}1) = \mathrm{val}(m)\cdot 2(m{+}1)/(2m{+}1)$.
\end{corollary}
\begin{proof}
Let $F(m) := {}_2F_1(\frac12, m{+}\frac12; m{+}\frac32; 1)$.
By Gauss summation,
$F(m)/F(m{+}1) = \Gamma(m{+}\frac32)\,\Gamma(m{+}2) / 
  (\Gamma(m{+}\frac52)\,\Gamma(m{+}1))
  = (m{+}1)/(m{+}\frac32) = 2(m{+}1)/(2m{+}3)$.
Therefore
\[
  \frac{\mathrm{val}(m{+}1)}{\mathrm{val}(m)}
  = \frac{2m{+}3}{2m{+}1}\cdot\frac{F(m)}{F(m{+}1)}
  = \frac{2m{+}3}{2m{+}1}\cdot\frac{2(m{+}1)}{2m{+}3}
  = \frac{2(m{+}1)}{2m{+}1}.  \qedhere
\]
\end{proof}

\begin{remark}[Wallis product interpretation]\label{rem:wallis-prod}
The shift relation shows $\mathrm{val}(m) = (2/\pi)\prod_{i=1}^m 2i/(2i{-}1)
= (2/\pi)\cdot 4^m/\binom{2m}{m}$:
the $S^{(m)}$ family consists of rational multiples of $2/\pi$,
where the rational coefficients are partial products of the
``rising half'' of the Wallis product for~$\pi/2$.
Equivalently, $1/\mathrm{val}(m)$ equals the normalised Wallis integral
$\int_0^{\pi/2}\!\sin^{2m}(t)\,dt \;/\; (\pi/2)
= \binom{2m}{m}/4^m$.
\end{remark}

\begin{corollary}[Full Pi Family]\label{cor:pi}
For all $m \ge 0$ the PCF with
$a(n) = -n(2n{-}(2m{+}1))$ and $b(n) = 3n+1$ converges to
$2^{2m+1}/(\pi\binom{2m}{m})$.
The family has been verified numerically to $600{+}$ stable digits
for $m = 0,\ldots,20$ at $1000$-digit working precision
(see Table~\ref{tab:1000dp}).
\end{corollary}

\begin{remark}[Status of former Conjecture]\label{rem:conj-resolved}
The binomial recurrence $\mathrm{val}(m{+}1) = \mathrm{val}(m)\cdot
2(m{+}1)/(2m{+}1)$, previously Conjecture~3.1 in earlier drafts,
is now proved as Corollary~\ref{cor:shift}
via the Gauss summation theorem for ${}_2F_1$ at unit argument.
\end{remark}

\begin{remark}[Parity Phenomenon]\label{rem:parity}
When $c = 2\ell$ is even, the numerator $a(n) = -n(2n{-}2\ell) = -2n(n{-}\ell)$
vanishes at $n = \ell$, so the continued fraction terminates
after $\ell$ steps and evaluates to a rational number.
For $c = 2,4,6,8,10,12,\ldots$ the values are
$1,\;3/2,\;15/8,\;35/16,\;315/128,\;693/256,\ldots$---the
Wallis product partial convergents.
\end{remark}

\begin{table}[h]
\centering
\caption{Pi Family: odd parameter $c = 2m{+}1$}\label{tab:pi}
\begin{tabular}{@{}ccccc@{}}
\toprule
$m$ & $c$ & $a(n)$ & $\mathrm{val}(m)\cdot\pi$ & $\mathrm{val}(m)$ \\
\midrule
$0$ & $1$ & $-n(2n{-}1)$ & $2$ & $2/\pi$ \\
$1$ & $3$ & $-n(2n{-}3)$ & $4$ & $4/\pi$ \\
$2$ & $5$ & $-n(2n{-}5)$ & $16/3$ & $16/(3\pi)$ \\
$3$ & $7$ & $-n(2n{-}7)$ & $32/5$ & $32/(5\pi)$ \\
$4$ & $9$ & $-n(2n{-}9)$ & $256/35$ & $256/(35\pi)$ \\
\bottomrule
\end{tabular}
\end{table}

\subsection{The Parity Bifurcation}\label{sec:bifurcation}

The parametric family $a(n) = -n(2n{-}c)$, $b(n) = 3n{+}1$
exhibits a sharp bifurcation depending on the parity of~$c$.

\subsubsection{The Rational Branch ($c \in 2\mathbb{Z}$)}\label{sec:rational-branch}

When $c = 2\ell$ is even, $a(\ell) = -\ell(2\ell - 2\ell) = 0$,
so the continued fraction terminates after $\ell$~steps.
The resulting values are Wallis product partial convergents:
$1,\;\frac{3}{2},\;\frac{15}{8},\;\frac{35}{16},\;\frac{315}{128},\;\frac{693}{256},\ldots$
(see Remark~\ref{rem:parity}).

\subsubsection{The Transcendental Branch ($c \in 2\mathbb{Z}+1$)}\label{sec:transcendental-branch}

When $c = 2m{+}1$ is odd, the PCF converges to
$\mathrm{val}(m) = 2^{2m+1}/(\pi\binom{2m}{m})$
(Corollary~\ref{cor:pi}).
Each $\mathrm{val}(m)$ is a \emph{rational multiple of $1/\pi$},
with the ``shift'' operator
\begin{equation}\label{eq:shift}
  \frac{\mathrm{val}(m{+}1)}{\mathrm{val}(m)}
  = \frac{2(m{+}1)}{2m{+}1}.
\end{equation}
This identity (Corollary~\ref{cor:shift}) has now been
proved via the Gauss summation theorem and independently verified
to $600{+}$ stable digits for all $m \in [0, 20]$
at~$1000$-digit working precision.
Table~\ref{tab:1000dp} presents the definitive
certification.

\begin{table}[h]
\centering
\caption{1000-digit verification of $\mathrm{val}(m)\cdot\pi = 2^{2m+1}/\binom{2m}{m}$}\label{tab:1000dp}
\begin{tabular}{@{}cclcc@{}}
\toprule
$m$ & $c$ & $\mathrm{val}(m)\cdot\pi$ & Stable digits & Ratio $\mathrm{val}(m)/\mathrm{val}(m{-}1)$ \\
\midrule
$0$ & $1$ & $2/1$                       & $604$ & --- \\
$1$ & $3$ & $4/1$                       & $613$ & $2/1$ \\
$2$ & $5$ & $16/3$                      & $621$ & $4/3$ \\
$3$ & $7$ & $32/5$                      & $629$ & $6/5$ \\
$4$ & $9$ & $256/35$                    & $636$ & $8/7$ \\
$5$ & $11$ & $512/63$                   & $643$ & $10/9$ \\
$6$ & $13$ & $2048/231$                 & $649$ & $12/11$ \\
$7$ & $15$ & $4096/429$                 & $655$ & $14/13$ \\
$8$ & $17$ & $65536/6435$               & $661$ & $16/15$ \\
$9$ & $19$ & $131072/12155$             & $667$ & $18/17$ \\
$10$ & $21$ & $524288/46189$            & $673$ & $20/19$ \\
$20$ & $41$ & $549755813888/34461632205$ & $725$ & $40/39$ \\
\bottomrule
\end{tabular}
\end{table}

\subsubsection{Convergence Analysis}\label{sec:conv-analysis}

The number of stable digits at depth~$N{=}2000$ increases with~$m$:
approximately $604 + 6.1m$ digits for $m \in [0,20]$
(a consequence of the slower growth of the error bound
$\prod_{j=1}^N |a(j)|/|b(j)|^2$ as $m$ increases).
At depth $N{=}4000$, all members $m = 0,\ldots,20$
achieve $600{+}$ stable digits at $1000$-digit working precision.
The family therefore admits \emph{arbitrary-precision certification}
for every member via standard backward-recurrence evaluation.

\section{Convergence and Bracketing}\label{sec:conv}

\begin{lemma}[Worpitzky Condition]\label{lem:work}
For both families, the normalised partial numerators
$|\tilde a_n| = |a_n|/(b_n\,b_{n-1})$ satisfy
$|\tilde a_n| < 1/4$ for all sufficiently large~$n$.
\end{lemma}
\begin{proof}
\textbf{Log family.}
The asymptotic ratio is $|\tilde a_n| \to k/(k{+}1)^2$.
Since $4k < (k{+}1)^2 \iff (k{-}1)^2 > 0$, we have
$k/(k{+}1)^2 < 1/4$ for all $k > 1$.
For $k \ge 2$, direct computation confirms $|\tilde a_n| < 1/4$
for all $n \ge 1$.
For $1 < k < 2$, the bound may fail at $n = 1$
(where $|\tilde a_1| = 1/(2k{+}1)$, which exceeds $1/4$
when $k < 3/2$); however, convergence for all $k > 1$ is
established independently by the closed-form proof
(Theorem~\ref{thm:log}).

\textbf{Pi family.}
$|\tilde a_n| = n(2n{-}1)/((3n{+}1)(3n{-}2)) \to 2/9 < 1/4$.
At $n = 1$: $|\tilde a_1| = 1\cdot1/(4\cdot1) = 1/4$, so strict inequality holds
for all $n \ge 2$.  By a classical result
(see~\cite[\S4.2]{lorentzen2008}), convergence still follows
when $|\tilde a_n| \le 1/4$ with equality at finitely many~$n$.
\end{proof}

\noindent Convergence of both families is therefore guaranteed:
for the log family by the closed-form proof when $1 < k < 2$ and
by Worpitzky for $k \ge 2$; for the pi family by Worpitzky
for all $n \ge 2$ (with the $n=1$ borderline handled by the
classical extension).

\begin{lemma}[Bracketing]\label{lem:bracket}
If $a_n < 0$ and $b_n > 0$ for all $n \ge 1$, and the CF converges,
then $\min(C_N,C_{N+1}) \le L \le \max(C_N,C_{N+1})$.
\end{lemma}
\begin{proof}
The determinant $p_nq_{n-1}-p_{n-1}q_n = \prod_{j=1}^n(-a_j)>0$
shows $C_n-C_{n-1}$ alternates in sign
(since $q_n>0$ by induction and $b_n>0$).
See~\cite[Theorem~4.35]{lorentzen2008}.
\end{proof}

\section{Hypergeometric Connections}\label{sec:hyper}

\textbf{Log family.}
From the proof of Theorem~\ref{thm:log}, the convergent ratio
$C_n = k/\sum_{j=0}^n (1/k)^j/(j{+}1)$ has limit determined by
the series $\sum_{j=0}^\infty x^j/(j{+}1) = {}_2F_1(1,1;2;x)$
evaluated at $x = 1/k$.  Since
$\ln(k/(k{-}1)) = (1/k)\cdot{}_2F_1(1,1;2;1/k)$, the PCF value
is $k/({}_2F_1(1,1;2;1/k))$.

\textbf{Pi family.}
The series~\eqref{eq:pi-id} arises from $\arcsin(1)=\pi/2$,
which equals ${}_2F_1(\frac12,\frac12;\frac32;1)$ via the Gauss summation theorem.
For general~$m$, Theorem~\ref{thm:pi-full} establishes that
\[
  \mathrm{val}(m)
  = \frac{2m{+}1}
         {{}_2F_1\!\bigl(\tfrac12,\,m{+}\tfrac12;\;m{+}\tfrac32;\;1\bigr)},
\]
where the Gauss evaluation
${}_2F_1(a,b;c;1) = \Gamma(c)\Gamma(c{-}a{-}b)/(\Gamma(c{-}a)\Gamma(c{-}b))$
(valid since $c - a - b = \frac12 > 0$) yields the closed form directly.
The parity phenomenon (Remark~\ref{rem:parity}) connects even-$c$
values to Wallis product partial products, while the odd-$c$
(transcendental) branch is governed by the Wallis integral
$\int_0^{\pi/2}\sin^{2m}(t)\,dt$
(Remark~\ref{rem:wallis-prod}).

\section{Numerical Certification}\label{sec:cert}

Arb ball arithmetic (\texttt{python-flint}~0.8.0, $10\,000$-bit precision)
at depth $N{=}5000$; by Lemma~\ref{lem:bracket}
the limit lies in $[\min(C_N,C_{N-1}),\max(C_N,C_{N-1})]$.

\begin{table}[h]
\centering
\caption{Arb-certified intervals (depth $5000$)}\label{tab:cert}
\begin{tabular}{@{}lccc@{}}
\toprule
PCF & Target & Bracket & Digits \\
\midrule
$a=-2n^2,\;b=3n{+}2$      & $1/\ln 2$    & $<10^{-1509}$ & $1509$ \\
$a=-3n^2,\;b=4n{+}3$      & $1/\ln(3/2)$ & $<10^{-2265}$ & $2265$ \\
$a=-n(2n{-}1),\;b=3n{+}1$ & $2/\pi$      & $<10^{-1508}$ & $1508$ \\
$a=-n(2n{-}3),\;b=3n{+}1$ & $4/\pi$      & $<10^{-1518}$ & $1518$ \\
$a=-n(2n{-}5),\;b=3n{+}1$ & $16/(3\pi)$  & $<10^{-621}$  & $621$ \\
$a=-n(2n{-}9),\;b=3n{+}1$ & $256/(35\pi)$& $<10^{-636}$  & $636$ \\
$a=-n(2n{-}21),\;b=3n{+}1$& $524288/(46189\pi)$ & $<10^{-673}$ & $673$ \\
\bottomrule
\end{tabular}
\end{table}

\section{Discovery Methodology}

\begin{enumerate}
\item \textbf{MITM sweep.}
  Fix $b(n)=dn{+}e$; enumerate $a(n)=pn^2{+}qn{+}r$
  with $|p|,|q|,|r|\le 5$; match against known constants.
\item \textbf{Modular signature sweep.}
  Vary $d\in\{2,\ldots,8\}$, $e\in\{1,\ldots,8\}$.
  The $d{=}3$ family produced $\pi$- and $\ln$-hits.
\item \textbf{Parametric generalization.}
  Sweep $c$ in $a(n)=-n(2n{-}c)$ and $k$ in $a(n)=-kn^2$.
\end{enumerate}

\section{Open Questions}

\begin{enumerate}
\item Is there a unified framework producing both families?
\item Theorem~\ref{thm:pi-m1} identifies $P_n^{(1)} = (2n{-}1)!!\,(n^2{+}3n{+}1)$.
  The denominator convergent $Q_n^{(1)}$ does not admit a polynomial--factorial
  closed form (consistent with $P_n/Q_n \to 4/\pi$ being transcendental).
  Can one find analogous formulas for $P_n^{(m)}$ when $m \ge 2$?
  (The closed form $\mathrm{val}(m) = 2^{2m+1}/(\pi\binom{2m}{m})$
  is now proved via the Gauss summation theorem in Theorem~\ref{thm:pi-full};
  an \emph{elementary} proof via direct induction on convergent sequences
  remains open for $m \ge 2$.)
\item Do these PCFs yield irrationality measures for $\ln 2$,
  $\ln(3/2)$, or $1/\pi$ via Ap\'ery-style arguments?
\item Do analogous families exist for $d \ge 5$,
  producing $\zeta(3)$, Catalan's constant, or polylogarithms?
\end{enumerate}

\section{Towards \texorpdfstring{$\zeta(3)$}{zeta(3)}: Quadratic GCF Candidates}\label{sec:zeta3}

The classical Ap\'ery CF for~$\zeta(3)$ uses sextic numerators:
$a(n)=-n^6$, $b(n)=34n^3-51n^2+27n-5$, converging to $6/\zeta(3)$.
A central open question is whether \emph{simpler} polynomial CFs
exist for~$\zeta(3)$ within the quadratic GCF family
$a(n) = c_0+c_1 n+c_2 n^2$, $b(n) = d_0+d_1 n+d_2 n^2$.

\begin{remark}[Classification of $\zeta(3)$ candidates]\label{rem:z3class}
Should a new CF converging to a rational multiple of~$\zeta(3)$
emerge from the quadratic GCF family, we propose the following
classification:
\begin{enumerate}
\item \textbf{Type~Q (quadratic--linear).}
  $\deg a = 2$, $\deg b = 1$.
  These are direct analogues of the pi family
  (Theorem~\ref{thm:pi}); a $\zeta(3)$ hit here would be
  the strongest possible result, since the convergents grow
  polynomially and the CF structure admits induction proofs.
\item \textbf{Type~C (cubic--quadratic).}
  $\deg a = 3$, $\deg b \le 2$.
  Ap\'ery-like cubic numerators $a(n) = -n^3 + \ldots$
  with quadratic~$b(n)$.  These are natural deformations
  of the $a(n)=-n^3$ seeds; closed forms may involve
  ${}_3F_2$ hypergeometric functions.
\item \textbf{Type~S (sextic--cubic).}
  $\deg a = 6$, $\deg b = 3$.
  Perturbations of Ap\'ery's original CF.
  Novel members of this family would extend the Ap\'ery
  phenomenon but likely resist elementary induction proofs.
\end{enumerate}
A discovery of Type~Q or~C with verified precision exceeding
$250$~digits and Wronskian stability (forward--backward
agreement) would constitute strong evidence for inclusion
in the Unified Discovery Theorem.
\end{remark}

\appendix
\section{Exact Convergents}\label{app:conv}

\textbf{Log ($k{=}2$):} $p_n = (n{+}1)!\cdot 2^{n+1}$.

\smallskip
\begin{center}
\begin{tabular}{@{}rrrl@{}}
\toprule
$n$ & $p_n$ & $q_n$ & $C_n$ \\
\midrule
0 & 2 & 1 & 2.000000 \\
1 & 8 & 5 & 1.600000 \\
2 & 48 & 32 & 1.500000 \\
3 & 384 & 262 & 1.465649 \\
4 & 3840 & 2644 & 1.452345 \\
5 & 46080 & 31832 & 1.447584 \\
\bottomrule
\end{tabular}
\end{center}
\smallskip\noindent Limit: $1/\ln 2 = 1.442695\ldots$

\medskip
\textbf{Pi ($m{=}0$):} $p_n = (2n{+}1)!!$.

\smallskip
\begin{center}
\begin{tabular}{@{}rrrl@{}}
\toprule
$n$ & $p_n$ & $q_n$ & $C_n$ \\
\midrule
0 & 1 & 1 & 1.000000 \\
1 & 3 & 4 & 0.750000 \\
2 & 15 & 22 & 0.681818 \\
3 & 105 & 160 & 0.656250 \\
4 & 945 & 1464 & 0.645492 \\
5 & 10395 & 16224 & 0.640688 \\
\bottomrule
\end{tabular}
\end{center}
\smallskip\noindent Limit: $2/\pi = 0.636620\ldots$

\section{Reproducibility}\label{app:repro}

Python~3.12, \texttt{mpmath}~1.4.1, \texttt{python-flint}~0.8.0 (Arb),
\texttt{sympy}~1.14.0.
Forward recurrence~\eqref{eq:rec} with $p_{-1}{=}1$, $p_0{=}b(0)$,
$q_{-1}{=}0$, $q_0{=}1$.

\begin{verbatim}
pip install mpmath python-flint sympy
python run_phase4_certification.py  # Arb
python run_phase5_proofs.py         # Proofs
python run_phase7_pi_analysis.py    # Pi convergents
\end{verbatim}

\noindent Runtime: ${<}\,2$~hours on standard desktop.
All source code, certification logs, and a verification notebook
are available at \url{https://github.com/PLACEHOLDER/pcf-families}
(to be made public upon publication).

\section{Multi-Precision Validation Toolkit}\label{app:toolkit}

For each candidate PCF identity $\PCF(a,b) \stackrel{?}{=} r\cdot\alpha$
(where $\alpha$ is a target constant such as $\zeta(3)$ and $r\in\mathbb{Q}$),
the following four-stage pipeline establishes confidence before
attempting a formal proof.

\begin{enumerate}
\item \textbf{Stage~1: PSLQ Matching} (working precision $P$).
  Evaluate the CF value $v$ at $P$ digits.
  Construct the vector $(v,\;\alpha,\;1)$ and run PSLQ
  (or LLL) to find integer relations $n_0 v + n_1\alpha + n_2 = 0$
  with bounded $|n_i| \le 100$.
  Accept if the residual $|n_0 v + n_1\alpha + n_2|$
  is smaller than $10^{-P+10}$.

\item \textbf{Stage~2: High-Precision Verification} (precision $\ge 2P$).
  Re-evaluate the CF and the target constant at $\ge 2P$ digits
  using \texttt{mpmath} or ball arithmetic (\texttt{python-flint}).
  Confirm that the PSLQ relation holds to $\ge 0.8 \times 2P$
  digits.  This eliminates coincidental matches whose probability
  decays as $10^{-P}$.

\item \textbf{Stage~3: Wronskian Stability Check.}
  Compare forward (Euler--Wallis) and backward (bottom-up)
  CF evaluations at $2P$ digits.
  Require $\ge 0.4 \times 2P$ digits of agreement.
  Disagreement signals ill-conditioning (small denominators,
  near-cancellation) that can produce spurious matches.

\item \textbf{Stage~4: Convergents Table.}
  Compute $C_N$ for $N \in \{10, 20, 50, 100, 200, 500, 1000\}$.
  Record $\Delta$-digits (agreement between consecutive depths).
  For a genuine identity, $\Delta$-digits should grow linearly
  with~$N$ (polynomial convergence) or faster.
  Stalling $\Delta$-digits indicate an accidental near-match
  that will eventually diverge.
\end{enumerate}

\begin{remark}[Irrationality Proof Strategy]
Once a candidate passes all four stages, a formal irrationality
proof proceeds as follows.
Given $\PCF(a,b) = r\cdot\alpha$ with $r = p/q \in \mathbb{Q}$:
\begin{itemize}
\item Extract convergent numerators~$P_n$ and denominators~$Q_n$
  from the recurrence~\eqref{eq:rec}.
\item Show that $|Q_n\alpha - P_n/r| = O(\lambda^{-n})$ for some
  $\lambda > 1$ (exponential convergence of the CF to its limit).
\item Verify that $Q_n \in \mathbb{Z}$ or that
  $\mathrm{lcm}(1,\ldots,n)^s Q_n \in \mathbb{Z}$
  for some fixed~$s$ (integrality condition).
\item If $\lambda > \mathrm{lcm}(1,\ldots,n)^s$ asymptotically,
  the standard Nesterenko criterion~\cite{nesterenko1996}
  yields irrationality (and an irrationality measure
  $\mu \le 1 + s\log n/\log\lambda$).
\end{itemize}
For the Ap\'ery CF ($a(n)=-n^6$, $b(n)=34n^3-51n^2+27n-5$),
$\lambda = (\sqrt{2}+1)^4$ and $s=3$ (accounting for
$\mathrm{lcm}(1,\ldots,n)^3$ clearing the denominators),
which gives Ap\'ery's theorem: $\zeta(3) \notin \mathbb{Q}$.
A Type~Q or Type~C candidate (Section~\ref{sec:zeta3}) with
lower polynomial degree would yield a \emph{better} irrationality
measure.
\end{remark}

\section*{Acknowledgements}

Inspired by the Ramanujan Machine~\cite{raayoni2021,elimelech2024}.

\begin{thebibliography}{5}

\bibitem{raayoni2021}
G.\,Raayoni et al.,
\emph{Generating conjectures on fundamental constants with the Ramanujan Machine},
Nature \textbf{590} (2021), 67--73.

\bibitem{elimelech2024}
D.\,Elimelech et al.,
\emph{Algorithm-assisted discovery of an intrinsic order among mathematical constants},
PNAS \textbf{121}(25) (2024), e2321440121.

\bibitem{cuyt2008}
A.\,Cuyt, V.\,B.\,Petersen, B.\,Verdonk, H.\,Waadeland, W.\,B.\,Jones,
\emph{Handbook of Continued Fractions for Special Functions},
Springer, 2008.

\bibitem{lorentzen2008}
L.\,Lorentzen, H.\,Waadeland,
\emph{Continued Fractions: Convergence Theory},
2nd ed., Atlantis Press, 2008.

\bibitem{osler2010}
T.\,J.\,Osler,
\emph{Lord Brouncker's forgotten sequence of continued fractions for pi},
Int.\ J.\ Math.\ Ed.\ Sci.\ Tech.\ \textbf{41} (2010), 249--256.

\bibitem{nesterenko1996}
Yu.\,V.\,Nesterenko,
\emph{On the linear independence of numbers},
Vestnik Moskov.\ Univ.\ Ser.~I Mat.\ Mekh.\ \textbf{1} (1985), 46--49;
English transl.\ Moscow Univ.\ Math.\ Bull.\ \textbf{40} (1985), 69--74.

\end{thebibliography}

\end{document}
